\chapter{Kidney Transplantation}
\section*{Overview}
Kidney transplantation places a healthy donor kidney into a recipient with end-stage kidney disease.
\section*{Common Instruments and Equipment}
The exact setup varies by technique and setting. Common items may include:
\begin{itemize}[nosep,leftmargin=*]
  \item vascular and transplant surgery instrument set
  \item vascular clamps and microsurgical sutures
  \item urinary catheter and ureteral stent when indicated
\end{itemize}

\section*{Why It Is Done}
It is considered for eligible patients with irreversible kidney failure as an alternative to long-term dialysis.
\section*{How It Is Performed}
The donor kidney is placed in the lower abdomen and connected to the recipient's blood vessels and bladder. The native kidneys usually remain unless they cause complications.
\section*{Contraindications}
No single contraindication applies to every patient; the treating team assesses whether the procedure is suitable for the individual and the reason for treatment.

\section*{Precautions}
The care team reviews relevant health conditions, medicines, allergies, and procedure-specific risks, and explains any changes needed for the individual.

\section*{Risks and Follow-up}
Rejection, infection, blood clots, urine leak, delayed graft function, medication toxicity, and recurrent disease are possible. Lifelong immunosuppression and transplant-team monitoring are required.
\section*{References}
\begin{enumerate}\item MedlinePlus. Kidney Transplant. U.S. National Library of Medicine; 2025.\item Current Medical Diagnosis and Treatment. McGraw-Hill; 2025.\end{enumerate}
\clearpage
